%% chapters/_template.tex
%% Plantilla para nuevos capítulos — Control Automático UTN IEE
%%
%% INSTRUCCIONES:
%%   1. Copiá este archivo a chapters/chapterXX_nombre.tex
%%   2. Reemplazá los marcadores [TEMA], [NÚMERO], etc.
%%   3. Descomentá las secciones que uses; eliminá las que no apliquen.
%%   4. Agregá \include{chapters/chapterXX_nombre} en main.tex

\chapter{[TÍTULO DEL CAPÍTULO]}
\label{ch:[etiqueta]}

%% Epígrafe opcional
\begin{flushright}
  \itshape
  ``[Cita relevante.]''\\[2pt]
  \normalfont\small --- [Autor], \textit{[Obra]}
\end{flushright}

\bigskip

%% Resumen introductorio del capítulo (3-5 líneas)
En este capítulo se estudia [breve descripción]. Al finalizar, el estudiante
será capaz de: (i)~[objetivo 1]; (ii)~[objetivo 2]; (iii)~[objetivo 3].

%======================================================================
\section{Introducción conceptual}
\label{sec:[etiqueta]-intro}
%======================================================================

[Texto introductorio que contextualiza el tema. Conectar con lo visto en
capítulos anteriores y motivar lo que viene.]

%======================================================================
\section{Motivación y aplicaciones reales}
\label{sec:[etiqueta]-aplicaciones}
%======================================================================

[Por qué este tema importa en la práctica. Ejemplos de ingeniería eléctrica.]

\begin{aplicacion}
[Descripción de una aplicación concreta en motores, redes eléctricas,
energías renovables, electrónica de potencia, etc.]
\end{aplicacion}

%======================================================================
\section{Fundamentos teóricos}
\label{sec:[etiqueta]-fundamentos}
%======================================================================

\begin{definicion}{[Nombre del concepto]}
  [Definición formal del concepto principal.]
\end{definicion}

[Desarrollo del marco teórico. Incluir intuición física antes de la matemática.]

\interfisica{[Qué significa físicamente este resultado.]}

%======================================================================
\section{Desarrollo matemático}
\label{sec:[etiqueta]-matematica}
%======================================================================

[Derivaciones paso a paso. Explicar cada línea. Usar entornos align para
ecuaciones de múltiples pasos.]

\begin{align}
  [ecuación 1] &= [resultado 1] \label{eq:[etiqueta]-1} \\
  [ecuación 2] &= [resultado 2] \label{eq:[etiqueta]-2}
\end{align}

\begin{teorema}{[Nombre del teorema]}
  [Enunciado formal. Hipótesis y tesis claramente indicadas.]
\end{teorema}

%======================================================================
\section{Interpretación física}
\label{sec:[etiqueta]-interpretacion}
%======================================================================

[Qué significa el resultado matemático en términos físicos. Usar analogías
cuando ayuden a la comprensión.]

\begin{observacion}
  [Observación relevante que el estudiante no debe pasar por alto.]
\end{observacion}

%======================================================================
\section{Ejemplos simples}
\label{sec:[etiqueta]-ejemplos-simples}
%======================================================================

\begin{ejemplo}{[Título del ejemplo]}
  \textbf{Enunciado.} [Descripción del problema.]

  \textbf{Solución.}
  \begin{enumerate}
    \item [Paso 1]
    \item [Paso 2]
  \end{enumerate}

  \textbf{Interpretación.} [Qué significa el resultado.]
\end{ejemplo}

%======================================================================
\section{Ejemplos avanzados}
\label{sec:[etiqueta]-ejemplos-avanzados}
%======================================================================

[Ejemplos más complejos que integran múltiples conceptos.]

% [INSERTAR FIGURA: descripción detallada del gráfico, qué muestra,
%  qué debe observar el alumno, por qué es pedagógicamente importante.]
\figplaceholder{[Descripción detallada del gráfico requerido.]}

%======================================================================
\section{Ejercicios resueltos}
\label{sec:[etiqueta]-ejercicios-resueltos}
%======================================================================

%% --- Ejercicio básico ---
\begin{ejercicioresuelto}{[Título] (Nivel básico)}
  \textbf{Enunciado.} [Descripción.]

  \textbf{Análisis conceptual.} [Qué tipo de problema es y qué herramientas aplican.]

  \textbf{Desarrollo matemático.}
  \begin{align*}
    [paso 1] &= [resultado] \\
    [paso 2] &= [resultado]
  \end{align*}

  \textbf{Interpretación del resultado.} [Qué dice el resultado.]
\end{ejercicioresuelto}

%% --- Ejercicio intermedio ---
\begin{ejercicioresuelto}{[Título] (Nivel intermedio)}
  [...]
\end{ejercicioresuelto}

%% --- Ejercicio tipo examen ---
\begin{ejercicioresuelto}{[Título] (Tipo examen)}
  [...]
\end{ejercicioresuelto}

%======================================================================
\section{Ejercicios propuestos}
\label{sec:[etiqueta]-ejercicios-propuestos}
%======================================================================

\begin{enumerate}
  \item [Enunciado ejercicio 1.]
  \item [Enunciado ejercicio 2.]
  \item [Enunciado ejercicio 3. \textit{Pista: [ayuda breve.]}]
  \item [Enunciado ejercicio 4.]
  \item [Enunciado ejercicio 5. Nivel avanzado.]
\end{enumerate}

%======================================================================
\section{Implementación en MATLAB}
\label{sec:[etiqueta]-matlab}
%======================================================================

[Introducción al código. Objetivo de la simulación.]

\begin{matlabcodetitulo}{[Descripción de lo que hace el script]}
% [Objetivo: qué se quiere demostrar]
% [Conexión con la teoría: qué concepto ilustra]

[código MATLAB]

% [Qué se espera observar en los gráficos resultantes]
\end{matlabcodetitulo}

[Explicación del código. Qué hace cada bloque. Qué mostrarán los gráficos.]

%======================================================================
\section{Gráficos sugeridos}
\label{sec:[etiqueta]-graficos}
%======================================================================

% [INSERTAR FIGURA: Gráfico 1 — descripción]
\figplaceholder{[Gráfico 1: descripción detallada.]}

% [INSERTAR FIGURA: Gráfico 2 — descripción]
\figplaceholder{[Gráfico 2: descripción detallada.]}

%======================================================================
\section{Resumen del capítulo}
\label{sec:[etiqueta]-resumen}
%======================================================================

En este capítulo se estudió [tema]. Los resultados más importantes son:

\begin{itemize}
  \item [Punto clave 1.]
  \item [Punto clave 2.]
  \item [Punto clave 3.]
\end{itemize}

%======================================================================
\section{Conceptos clave}
\label{sec:[etiqueta]-conceptos}
%======================================================================

\begin{conceptosclave}
\begin{description}[leftmargin=4em, labelwidth=3.5em, font=\bfseries\color{colkey}]
  \item[Término 1]   Definición breve y precisa.
  \item[Término 2]   Definición breve y precisa.
  \item[Término 3]   Definición breve y precisa.
\end{description}
\end{conceptosclave}

%======================================================================
\section{Errores comunes}
\label{sec:[etiqueta]-errores}
%======================================================================

\begin{errorcomon}[Error 1]
  [Descripción del error. Por qué ocurre y cómo evitarlo.]
\end{errorcomon}

\begin{errorcomon}[Error 2]
  [Descripción del error. Por qué ocurre y cómo evitarlo.]
\end{errorcomon}
